% ЗГЕНЕРОВАНО scripts/build_format_v2_chapter.py — не редагувати вручну.
% Джерело: books/ashtavakra-gita/poetychnyi/rozdil-14.md

\poemverse{14.1}{Хто за природою має порожній розум — \\ навіть якщо з недбальства плекає в собі стани, — \\ той наче заснув і прокинувся: \\ блукання його вичерпане.}
\poemverse{14.2}{Де багатства, де друзі, \\ де мої слуги — чуттєві предмети? \\ Де писання, де знання — \\ коли в мене розтанула спрага?}
\poemverse{14.3}{Коли пізнано пурушу-свідка, \\ вищого Атмана й Ішвару — \\ а також безнадію, узу й визволення — \\ у мене немає турботи про визволення.}
\poemverse{14.4}{У того, хто всередині порожній від розрізнень, \\ а зовні ходить, як сам захоче, \\ стани бувають усякі, наче в божевільного; \\ і знають його тільки такі самі, як він.}
